\scanpage{74}
% archive foliation: 73
coïncident dans les $T_i\cap T_j$, donc définissent
$T\to X$ et ce dernier satisfait aux conditions voulues.
D'ailleurs, \deleted{\unsure{tt}} solution $u$ est, à priori, obtenue
par une construction comme celle qui précède.
Introduisant alors les \unsure{pour} $u,u'$ \unsure{liés} à :
$(T_i)$, $(T'_j)$, le recouvrement commun $(T_i\cap T'_j)_{i,j\in I\times I'}$,
on déduit l'unicité.

\uline{Lemme 3} La condition $\beta$ est automatiquement vérifiée
\inserted{(pour tt rec. $\mathfrak X=\bigcup\mathfrak X_i$ et $\mathfrak X$ quelconque)}
si on a un isomorphisme de foncteurs
\[
|T|\xrightarrow{\sim}|\underline W(T)|
\]
en particulier dans les cas $\mathrm{II}$, \deleted{$\unsure{\mathrm{I}}$},
en particulier $\mathrm{III}$.

C'est trivial.

\uline{Lemme 4} La condition $\alpha$ est automatiquement
vérifiée dans les cas $\mathrm{I}'$, $\unsure{\mathrm{II}^{a}|\mathrm{II}'}$.

C'est trivial pour $\unsure{\mathrm{II}^{a}|\mathrm{II}'}$.
Pour $\mathrm{I}'$, petit raisonnement, utilisant qu'un morphisme
fini est \uline{fermé}. (\unsure{la hyp. utile})
